\documentclass[11pt, oneside]{article}   	% use "amsart" instead of "article" for AMSLaTeX format
\usepackage{geometry}                		% See geometry.pdf to learn the layout options. There are lots.
\geometry{letterpaper}                   		% ... or a4paper or a5paper or ... 
%\geometry{landscape}                		% Activate for rotated page geometry
\usepackage[parfill]{parskip}    		% Activate to begin paragraphs with an empty line\usepackage{graphicx}				% Use pdf, png, jpg, or eps§ with pdflatex; use eps in DVI mode
								% TeX will automatically convert eps --> pdf in pdflatex
								
%\usepackage{{Gratzer-Color-Scheme}	% Active to color theorems red and lemmas blue
	
\usepackage{amsmath}
\usepackage{amssymb}

%SetFonts

%SetFonts


\title{Problem Set 1}
\author{Brian Ledger}
%\date{}							% Activate to display a given date or no date

\begin{document}
\maketitle
\section*{Problem 1}
\emph{A random variable $X$ takes values $\{1,2,4\}$ with probabilities $\{0.2,0.5,0.3\}$:}

$$
P(X=1) = 0.2, \quad P(X=2) = 0.5, \quad P(X = 4) = 0.3
$$
\subsection*{a)}
\emph{Compute the expected value $\mathbb{E}[X]$.}

The formula for expected value is $\mathbb{E}[X] = \sum^{x_i} x_iP(X=x_i)$

The expected value for X is

$$
\mathbb{E}[X] = \sum^{x_i} x_iP(X=x_i) = 1*0.2 + 2 * 0.5 + 4 * 0.3 = 2.4
$$

\subsection*{b)}
\emph{Explain what the expected value represents in this context.}

The expected value is the weighted sum of all possible outcomes, where the weight of each outcome is its respective probability.

\section*{Problem 2}
\emph{Using the random variable X from Problem 1:}

\subsection*{a)}
\emph{Compute $\mathbb{E}[X^2]$.}

$$
\mathbb{E}[X^2] = \sum^{x_i} x_i^2P(X=x_i) = 1*0.2 + 4 * 0.5 + 16 * 0.3 = 8.8
$$

\subsection*{b)}
\emph{Compute the variance $Var(X)$.}

The formula for variance is $Var(X) = \mathbb{E}[(X-\bar{X})^2] = \mathbb{E}[X^2]-\mathbb{E}[X]^2$.

$$
\mathbb{E}[X^2]-\mathbb{E}[X]^2 = 8.8 - 2.4^2 = 3.04
$$

\subsection*{c)}
\emph{Compute the standard deviation $\sigma_X$.}

The mean of $X$ is $1 + 2 + 4 / 3 = 7/3 = 2.\bar{3}$.

The formula for standard deviation of a sample is:

$$
\sigma = \sqrt{\frac{\sum^{x_i} (x_i-\bar{X})^2}{N-1}}
$$

\begin{eqnarray*}
\sigma &=& \sqrt{\frac{\sum^{x_i} (x_i-\bar{X})^2}{N-1}}\\
&=& \sqrt{\frac{(1-2.\bar{3})^2 + (2-2.\bar{3})^2 + (4-2.\bar{3})^2}{2}}\\
&=& \sqrt{\frac{-1.\bar{3}^2 + -0.\bar{3}^2 + 1.\bar{6}^2}{2}}\\
&=& \sqrt{\frac{-1.\bar{7} + -0.\bar{1} + 2.\bar{7}}{2}}\\
&=& \sqrt{\frac{-1.\bar{7} + -0.\bar{1} + 2.\bar{7}}{2}}\\
&=& \sqrt{\frac{4.\bar{6}}{2}}\\
&=& \sqrt{2.\bar{3}}\\
&\approx&1.528
\end{eqnarray*}


\subsection*{d)}
\emph{Interpret the meaning of the standard deviation.}

The standard deviation measures dispersion about a mean.

Approximately 68\% of data fall within 1.528 of the mean.

\section*{Problem 3}
\subsection*{a)}
\emph{A trader receives \$1 if a stock goes up tomorrow and \$0 otherwise. The probability the stock goes up is $p = 0.6.$}

\subsection*{b)}
\emph{Define the random variable X.}

$$
P(X) = 
 \begin{cases} 
      0.4, & X = \$0 \\
      0.6, & X = \$1
   \end{cases}
$$

\subsection*{c)}
\emph{Identify the distribution of X.}

The distribution of $X$ is discrete, binary. This is a binomial distribution. A binomial distribution with one trial is a Bernoulli distribution.

\subsection*{d)}
\emph{Compute $\mathbb{E}[X]$, $Var(X)$, and $\sigma_{X}$.}

$$
\mathbb{E}[X] = \$0*0.4 + \$1*0.6 = \$0.60
$$

$$
\bar{X} = \frac{0+1}{2} = 0.5
$$

$$
Var_{population}(X) = \frac{-0.5^2 + 0.5^2}{2} = 0.25
$$

$$
Var_{sample}(X) = \frac{-0.5^2 + 0.5^2}{1} = 0.5
$$

$$
\sigma^{population}_X = \sqrt{Var_{population}(X)} = 0.5
$$

$$
\sigma^{sample}_X = \sqrt{Var_{sample}(X)} \approx 0.707
$$

\section*{Problem 4}
\emph{A salesperson makes n=5 independent sales calls, each with probability of success $p = 0.3$.}

\subsection*{a)}
Define the random variable X.

$$
X \sim B(n=5, p = 0.3)
$$

$X$ is the number of successful sales calls in a trial of 5 independent sales calls. The distribution of $X$ successes in $n$ independent trials, with probability of success, $p$, is a Binomial Distribution.

\subsection*{b)}
\emph{Identify the distribution of X.}

The general formula for $k$ successes in $n$ independent trials, with probability of success, $p$ is:

$$
Pr(X = k) = \binom{n}{k} p^k*(1-p)^{n-k}
$$

With $p=0.3$, this evolves into:

\begin{eqnarray*}
Pr(X=k)&=&\binom{n}{k} 0.3^k*0.7^n*0.7^{-k}\\
&=& \binom{n}{k} 0.7^n*e^{-k(ln(0.7-ln(0.3))}\\
&=& \binom{n}{k} 0.7^n*e^{-0.847k}
\end{eqnarray*}

With $n=5$, this evolves into:
\begin{eqnarray*}
\binom{5}{k} 0.7^5*e^{-0.847k}
&=& \binom{5}{k} 0.168e^{-0.847k}\\
&=& \frac{5!*0.168e^{-0.847k}}{k!(5-k)!}\\
&=& \frac{120 * 0.168e^{-0.847k}}{k!(5-k)!}\\
&=& \frac{20.16e^{-0.847k}}{k!(5-k)!}
\end{eqnarray*}

Tabulated over the 6 possible outcomes, with 3 significant figures:

\begin{tabular}{|l|l|}\hline
k&P(X=k)\\\hline
0&0.168\\\hline
1&0.360\\\hline
2&0.308\\\hline
3&0.132\\\hline
4&0.028\\\hline
5&0.002\\\hline
\end{tabular}

\subsection*{c)}
\emph{Compute $P(X = 3)$.}

$$
P(X = 3) \approx 0.132
$$

\subsection*{d)}
\emph{Compute $\mathbb{E}[X]$ and $Var(X)$.}

The expectation for a binomial distribution is $\mathbb{E}[x] = np$. With $n=5$ and $p=0.3$), we have $\mathbb{E}[x] = 1.5.$

The variance for a binomial distribution is $Var(X) = np(1-p)$. With $n=5$ and $p=0.3$), we have $Var(X) = 1.05$.

\section*{Problem 5}
\emph{Using the binomial model from Problem 4:}

\subsection*{a)}
\emph{What does the standard deviation tell you about the number of successful calls?}

The standard deviation measures the expected dispersion around the mean.

In this case, the s.d. is ~1.025, so you may expect 1 to 2 out of 5 calls to succeed, most of the time.

\subsection*{b)}
\emph{Would you expect most outcomes to be close to the expected value? Why?}

The expected value is the projected average over many calls. The standard deviation of ~1.025 around the expected value of 1.5 indicates that most trials will see 2-3 successes. This seems likely, given that 2 and 3 successes are the most probable of outcomes in a trial of 5 calls.

\section*{Problem 6}
\emph{The number of customer complains arriving per hour follows a Poisson distribution with mean $\lambda = 2$ complaints/hour.}
\subsection*{a)}
\emph{Write the probability mass function.}

\subsection*{b)}
\emph{Compute $P(X=3)$.}

\subsection*{c)}
\emph{Compute $\mathbb{E}[X]$, $Var(X)$, and $\sigma_X$.}

\subsection*{d)}
\emph{Explain why the Poisson distribution is appropriate here.}

\section*{Problem 7}
\emph{Compare the variance and expected value of a Poisson distribution.}

\subsection{a)}
\emph{What special property do you observe?}

\subsection{b)}
\emph{How does this differ from the binomial distribution?}

\section*{Problem 8}
\emph{Answer briefly:}

\subsection{a)}
\emph{When should a Bernoulli distribution be used instead of a binomial?}

\subsection{b)}
\emph{Why does the Poisson distribution not require a fixed number of trials?}

\subsection{c)}
\emph{In practical terms, what does a high standard deviation imply?}

\end{document}  